% Plan: development/outline.md, section 5, "§4". Proofs of the validity
% results and of thm:lnts-opt (with lem:lnts-elim, prop:lnts-support):
% Appendix G-proofs-split. Full proofs, computations and verification
% records of the certificates: supplement files B1 (lnts, lukvle10),
% B2 (dtoc5, optcdeg2), B4 (chain, catmix) and B8 (waterno2). The certificate
% boxes live in those supplement sections (round-1 adjudication G2-09); here one
% sentence per theorem names arithmetic, implementations and evidence level.
% Replay tiers are only in Table 7 and the boxes (round-2 adjudication G2-02).
% Part (c) of prop:split-affine and the last statement of prop:split-cellwise,
% with their proofs, moved to the supplement (prop:split-forced,
% prop:split-cellwise-best in B10-split-extras.tex; round-2 G2-01).
\section{Split certificates along stage structure}\label{sec:split}

Fifteen of the 31 closures form the staged class of \cref{tab:closures}: \inst{lnts50}--\inst{lnts400}, \inst{dtoc5}, \inst{optcdeg2}, \inst{lukvle10}, \inst{chain50}--\inst{chain400} and \inst{catmix100}--\inst{catmix800}; the \inst{waterno2} bounds use the same idea.
We cut the model along its stages, add to each stage a function of the variables that it shares with the next stage, and subtract the same function from the next stage, so that these terms cancel at every feasible point; rigorous lower bounds on the stage minima then sum to a dual bound.
The \inst{catmix} certificates use a variant: they bound the value functions of the stage recursion from below, and their validity rests on concavity of the true value functions (\cref{lem:split-minorant}).
\Cref{app:splitproofs} proves the validity results of \cref{sec:split-framework}; \cref{tab:staged} summarizes the certificates, and \cref{app:splitproofs-lnts,app:families} give their full proofs, certificate boxes and verification records.

\subsection{Path model and validity}\label{sec:split-framework}

Let $N\ge1$.
Bag $t\in\{1,\dots,N\}$ (a block of variables, as in a tree decomposition) holds a vector $z_t\in\R^{d_t}$, a set $K_t\subseteq\R^{d_t}$ (the bounds and the rows that involve only bag~$t$) and a cost $F_t\colon\R^{d_t}\to\R\cup\{+\infty\}$.
Consecutive bags share a separator: $\pi^R_tz_t=\pi^L_{t+1}z_{t+1}\in\R^{k_t}$ for $1\le t<N$, where $\pi^R_t$ and $\pi^L_{t+1}$ select coordinates.
The path model is
\begin{equation}\label{eq:split-path}
  \text{(P)}\qquad
  \vstar=\inf\Bigl\{\sum_{t=1}^NF_t(z_t) : z_t\in K_t,\ F_t(z_t)<\infty\ (1\le t\le N),\ \pi^R_tz_t=\pi^L_{t+1}z_{t+1}\ (1\le t<N)\Bigr\},
\end{equation}
with $\vstar=+\infty$ if its feasible set $\feas[\mathrm P]$ is empty.
The costs $F_t$ take the value $+\infty$ where an expression of the model is undefined; this encodes the domain rule of \cref{def:sem-model}(iv).
For a transcribed control problem, bag~$t$ is the stage $(x_t,u_t,x_{t+1})$ with its dynamics rows, and its separator is $x_{t+1}$; a variable of every stage, such as the step $h$ of \inst{lnts}, lies in every separator.

A \emph{split} is a tuple $\varphi=(\varphi_1,\dots,\varphi_{N-1})$ of real functions $\varphi_t\colon\R^{k_t}\to\R$; we put $\varphi_0=\varphi_N=0$.
For sets $K'_t\subseteq\R^{d_t}$ define
\begin{equation}\label{eq:split-bound}
  F^\varphi_t(z)=F_t(z)+\varphi_t(\pi^R_tz)-\varphi_{t-1}(\pi^L_tz),
  \qquad
  B(\varphi;K')=\sum_{t=1}^N\inf_{z\in K'_t}F^\varphi_t(z),
\end{equation}
with $\inf\emptyset=+\infty$ and the convention that a sum with a term $+\infty$ equals $+\infty$; write $B(\varphi)=B(\varphi;K)$.
We call $F^\varphi_t$ the \emph{stage residual} of bag~$t$.
We also call bag~$t$ the $t$-th stage; its \emph{stage problem} is the minimization of $F^\varphi_t$ over $K'_t$, with value the \emph{stage infimum}.
We call $K'_t$ an \emph{enclosure} if $K'_t\supseteq\proj_t\feas[\mathrm P]$, the set of bag-$t$ components of feasible points.
The \emph{split form} is the set of functions allowed for the $\varphi_t$: affine, quadratic, cellwise affine, and so on.

\begin{lemma}[split bound]\label{lem:split-bound}
Let $\varphi$ be a split.
\begin{enumerate}
\item[(a)] If every $K'_t$ is an enclosure, then $B(\varphi;K')\le\vstar$.
\item[(b)] \emph{Infeasibility.} Let every $K'_t$ be an enclosure.
  If $B(\varphi;K')=+\infty$, or if every $F_t$ vanishes on $K_t\cap\operatorname{dom}F_t$ and $B(\varphi;K')>0$, then $\feas[\mathrm P]=\emptyset$.
\item[(c)] \emph{Case split.} If $\feas[\mathrm P]\subseteq Z_1\cup\dots\cup Z_m$ and $\beta_c\le\inf\{\sum_tF_t(z_t):z\in\feas[\mathrm P]\cap Z_c\}$ for each $c$, then $\min_c\beta_c\le\vstar$.
  For a box $Z_c$, $\beta_c$ may be a bound of form (a) for the model restricted to $Z_c$, with its own split and enclosures.
\item[(d)] \emph{Row form.} Let $\psi_1,\dots,\psi_r$ be real functions of $z=(z_1,\dots,z_N)$ that vanish on $\feas[\mathrm P]$, such as linear combinations of rows, and let $\mu\in\R^r$ and $Y\supseteq\feas[\mathrm P]$.
  Then $\inf_{z\in Y}\bigl[\sum_tF_t(z_t)+\sum_k\mu_k\psi_k(z)\bigr]\le\vstar$, and $\feas[\mathrm P]=\emptyset$ if this infimum is positive and every $F_t$ vanishes on $K_t\cap\operatorname{dom}F_t$.
  With the copy rows $\pi^R_tz_t-\pi^L_{t+1}z_{t+1}$ as the $\psi_k$ and $Y=K'_1\times\dots\times K'_N$, this bound is $B(\varphi;K')$ for an affine split (\cref{prop:split-affine}).
\end{enumerate}
\end{lemma}

Since $\sum_tF^\varphi_t(z_t)=\sum_tF_t(z_t)$ at every feasible point, a certificate of form (a) must establish a rigorous lower bound on every stage infimum and the validity of the enclosures.

For nonempty control sets $U_i$ and maps $g_i\colon Y_{i-1}\times U_i\to Y_i$ ($1\le i\le n$), consider the dynamics $y_i=g_i(y_{i-1},u_i)$ with $y_i\in Y_i$ and $u_i\in U_i$.
For a terminal cost $\Phi\colon Y_n\to\R$, let $V_n=\Phi$ and $V_{i-1}(y)=\inf_{u\in U_i}V_i(g_i(y,u))$, so that $\vstar=\inf_{Y_0}V_0$.

\begin{lemma}[value-function minorants]\label{lem:split-minorant}
\begin{enumerate}
\item[(a)] If $W_0\le V_0$ on $Y_0$, then $\inf_{Y_0}W_0\le\vstar$.
\item[(b)] \emph{Chord interpolation.} Let $Y_{i-1}=Y_i=\R^2_+$ and $g_i(y,u)=M_i(u)y$ with every $M_i(u)$ entrywise nonnegative, and let $V_{i-1}$ be real-valued, concave and positively homogeneous on $\R^2_+$.
  Let $W_i\le V_i$ on $\R^2_+$, let $r_0=(1,0),r_1,\dots,r_J=(0,1)$ be vectors in $\R^2_+$ with strictly increasing angles, and let $w_j\le\inf_{u\in U_i}W_i(M_i(u)r_j)$ for every $j$.
  Then the chord interpolant, $W_{i-1}(\sigma r_j+\sigma'r_{j+1})=\sigma w_j+\sigma'w_{j+1}$ for $\sigma,\sigma'\ge0$, satisfies $W_{i-1}\le V_{i-1}$ on $\R^2_+$.
\end{enumerate}
\end{lemma}

Part~(b) uses concavity of the true $V_{i-1}$, which lies above its chords between rays, not of the computed $W_i$, so \cref{lem:split-minorant} is not an instance of \cref{lem:split-bound}.

\begin{proposition}[affine splits]\label{prop:split-affine}
Let $\varphi_t(s)=\lambda_t^\top s+c_t$ with $\lambda_t\in\R^{k_t}$ and $c_t\in\R$, and put $\lambda_0=\lambda_N=0$.
\begin{enumerate}
\item[(a)] $B(\varphi)$ does not depend on $c$.
  It is the Lagrangian dual function, at $\lambda$, of the copy formulation, in which every bag has its own copy of each separator coordinate and the copy rows $\pi^R_tz_t-\pi^L_{t+1}z_{t+1}=0$ are dualized with multipliers $\lambda_t$.
\item[(b)] \emph{Exactness.} If some $\bar z\in\feas[\mathrm P]$ has $\bar z_t\in\argmin_{K_t}F^\varphi_t$ for every $t$, then $B(\varphi)=\vstar=\sum_tF_t(\bar z_t)$, and $\bar z$ is a global minimizer.
  Conversely, if $B(\varphi)=\vstar$ and $\vstar$ is attained at $z^*$, then $z^*_t\in\argmin_{K_t}F^\varphi_t$ for every $t$.
\end{enumerate}
\end{proposition}

The slopes of an exact affine split are determined by the model: at a minimizer whose bag components are interior points of the $K_t$ at which the costs $F_t$ are differentiable, they are the copy-row multipliers, and after the dynamics rows are eliminated they are the discrete costates (\cref{prop:split-forced}, explanatory).
The affine certificates below take their slopes from a local KKT point; validity never depends on this choice.
Part~(b) is a stagewise exactness test; for convex stage problems at a stationary point it is a discrete form of the sufficiency condition of \citet{mangasarian1966-sufficient-conditions-for-the-optimal}.

\begin{proposition}[windows]\label{prop:split-window}
Let $K'_t$ be enclosures, $1\le a\le b\le N$, and
\begin{multline*}
  \beta_{[a,b]}(\varphi;K')=\inf\Bigl\{\sum_{t=a}^bF_t(z_t)+\varphi_b(\pi^R_bz_b)-\varphi_{a-1}(\pi^L_az_a) :\\
  z_t\in K'_t\ (a\le t\le b),\ \pi^R_tz_t=\pi^L_{t+1}z_{t+1}\ (a\le t<b)\Bigr\}.
\end{multline*}
Then $B(\varphi;K')\le\sum_{t\notin[a,b]}\inf_{K'_t}F^\varphi_t+\beta_{[a,b]}(\varphi;K')\le\vstar$.
Disjoint windows combine likewise.
\end{proposition}

A window loses nothing on its interior separators, at the price of one global problem in its free variables.

\begin{proposition}[cellwise slopes]\label{prop:split-cellwise}
For each separator $t\in\{1,\dots,N-1\}$ let $\mathcal P_t$ be a finite family of sets (cells) whose union contains $\pi^R_tz_t$ for every $z\in\feas[\mathrm P]$, and give each cell $D\in\mathcal P_t$ a slope $\lambda_{t,D}\in\R^{k_t}$.
Put $\mathcal P_0=\mathcal P_N=\{\ast\}$ with zero slope, and read the conditions $\pi^L_1z\in\ast$ and $\pi^R_Nz\in\ast$ as void.
Let $K'_t$ be enclosures, and for every $t$ and $(D,D')\in\mathcal P_{t-1}\times\mathcal P_t$ let $b_t(D,D')\in\R\cup\{+\infty\}$ satisfy
\[
  b_t(D,D')\le\inf\bigl\{F_t(z)+\lambda_{t,D'}^\top\pi^R_tz-\lambda_{t-1,D}^\top\pi^L_tz :
  z\in K'_t,\ \pi^L_tz\in D,\ \pi^R_tz\in D'\bigr\}.
\]
Then
\[
  \vstar\ \ge\ \mathrm{SP}:=\min\Bigl\{\sum_{t=1}^Nb_t(D_{t-1},D_t) : D_t\in\mathcal P_t\ (1\le t<N),\ D_0=D_N=\ast\Bigr\},
\]
the shortest-path value in the layered graph with arc weights $b_t$.
\end{proposition}

The only coupling condition is that the two bags at a separator use the same slope for the same cell; constants per cell cancel along every path, so a certificate stores slopes only.
With exact pair bounds on a partition and finite $\mathrm{SP}$, $\mathrm{SP}$ is also the largest bound that cellwise-affine splits with these slopes give (\cref{prop:split-cellwise-best}, explanatory).

\begin{remark}[derived enclosures]\label{rem:split-enclosures}
Several certificates use derived enclosures of unbounded states: the \inst{optcdeg2} velocities (\cref{lem:optcdeg2-enclosure}), the \inst{chain} end values (\cref{lem:chain-window}), the \inst{lukvle10} stage minimizers (\cref{lem:lukvle10-coercive}) and the \inst{waterno2} tank levels (\cref{lem:waterno2-levels}).
\end{remark}

\Cref{lem:split-bound}(a) is Krotov's sufficient condition for discrete control systems \citep{krotov1967-sufficient-conditions-for-the-optimality} as a Bellman inequality, related to relaxed dynamic programming \citep{lincoln2006-relaxing-dynamic-programming}: the split functions are Krotov's verification functions, and a split with $B(\varphi)=\vstar$ is a discrete calibration in the sense of the calculus of variations.
\Cref{prop:split-affine} is Lagrangian decomposition of copy rows, as in generalized Benders decomposition \citep{geoffrion1972-generalized-benders-decomposition} and in Lagrangian bounds within spatial branch and bound \citep{karuppiah2008-a-lagrangean-based-branch-and,khajavirad2009-a-deterministic-lagrangian-based-global,cao2019-a-scalable-global-optimization-algorithm}; and \cref{prop:split-cellwise} is closest to multistage decomposition with state copies and Lagrangian cuts \citep{zhang2022-stochastic-dual-dynamic-programming-for,fullner2022-non-convex-nested-benders-decomposition}, in particular cuts that differ across a partition of the state range \citep{yang2025-globally-converging-algorithm-for-multistage}.

\subsection{Exact affine splits: \inst{lnts} and \inst{dtoc5}}\label{sec:split-affine}

For $N\in\{50,100,200,400\}$, \inst{lnts}$N$ is the trapezoidal transcription with $N$ steps of length $h$ of the particle-steering problem of COPS~2.0 \citep{dolan2001-benchmarking-optimization-software-with-cops}.
A thrust of magnitude 100 in a steerable direction $\theta$ accelerates a particle from rest at the origin to height~5 and velocity $(45,0)$ in minimum time $Nh$.
With angles $\theta_j$, positions $p_j=(p^x_j,p^y_j)$, velocities $v_j=(v^x_j,v^y_j)$ ($0\le j\le N$) and $e(\theta)=(\cos\theta,\sin\theta)$, the model minimizes $Nh$ subject to
\[
  v_{i+1}-v_i=50h\bigl(e(\theta_i)+e(\theta_{i+1})\bigr),\qquad p_{i+1}-p_i=\tfrac h2(v_i+v_{i+1})\qquad(0\le i<N),
\]
$h\ge0$, zero initial states, $p^y_N=5$, $v^x_N=45$, $v^y_N=0$ and $|\theta_j|\le\bar\theta:=1.5707963267949$, a decimal that exceeds $\pi/2$ by between \sci{3}{-15} and \sci{4}{-15}; all other states are free.
Let $w_0=w_N=\frac12$ and $w_j=1$ otherwise; $c_0=(2N-1)/4$, $c_j=N-j$ for $0<j<N$ and $c_N=\frac14$; $r_j=c_j/w_j-N/2$; and $\alpha(h)=9/(20h)$, $\beta(h)=1/(20h^2)$.
For $\nu\ge0$ define
\begin{equation}\label{eq:lnts-g}
  C(\nu)=\sum_{j=0}^N\frac{w_j}{\sqrt{1+\nu^2r_j^2}},\qquad
  D(\nu)=\nu\sum_{j=0}^N\frac{w_jr_j^2}{\sqrt{1+\nu^2r_j^2}},\qquad
  g(\nu)=D(\nu)-\frac{20}{81}C(\nu)^2.
\end{equation}

\begin{theorem}[\inst{lnts} exact optimum]\label{thm:lnts-opt}
Let $N\in\{50,100,200,400\}$.
The function $g$ has exactly one positive root $\nu^*$.
The optimal value of \inst{lnts}$N$ is $\vstar=Nh^*$ with $h^*=9/(20\,C(\nu^*))$.
It is attained by $\theta^*_j=\arctan(\nu^*r_j)$, $h=h^*$ and the states that the dynamics determine.
The optimal values satisfy
\[
\begin{array}{ll}
  N=50: & 0.5546687649386788\le\vstar\le0.5546687649386789,\\
  N=100: & 0.5545954011669111\le\vstar\le0.5545954011669112,\\
  N=200: & 0.5545770161030836\le\vstar\le0.5545770161030837,\\
  N=400: & 0.5545724137006871\le\vstar\le0.5545724137006872.
\end{array}
\]
\end{theorem}

\Cref{app:splitproofs-lnts} gives the proof.
Summing the dynamics rows reduces exact feasibility to $h>0$, the angle bounds and three identities in the controls (\cref{lem:lnts-elim}).
At a fixed step $h$ the objective is the constant $Nh$; dropping it, combining these identities with multipliers $(\mu,\nu)$ and applying the Cauchy--Schwarz inequality to each control gives a support bound; this is \cref{lem:split-bound}(d) in its infeasibility form (\cref{prop:lnts-support}).
At $\mu=-\nu^*N/2$ the controls $\theta^*_j$, which follow the discrete linear-tangent law $\tan\theta^*_j=\nu^*(N/2-j)$ for $0<j<N$, attain the support bound; this proves optimality, and rational brackets of $\nu^*$ of width below $10^{-90}$ give the enclosures.
In the certificate (\cref{app:lnts-proof}), two separately written integer-only codes enclose the optima in arithmetic \arith{E}, one of them to the displayed digits; its evidence level is \evid{hand}+\evid{stored}.

\inst{dtoc5} is a discrete-time control problem with $T=49{,}999$ steps and $h=1/50000$:
\[
  \min\ h\sum_{t=0}^{T-1}(u_t^2+y_t^2)\quad\text{s.t.}\quad y_{t+1}=y_t+4hy_t^2-hu_t\quad(0\le t<T),\qquad y_0=1,
\]
with all other variables free; the final state $y_T$ does not enter the objective.
The CUTEst problem DTOC5 has $h$ in place of $4h$ \citep{luksan2026-cutest-source-forms-for-dtoc5,coleman1993-the-global-convergence-of-a}; our results concern the MINLPLib model only.

\begin{proposition}[\inst{dtoc5} dual identity]\label{prop:dtoc5-identity}
Write $c_t(x)=y_{t+1}-y_t-4hy_t^2+hu_t$, the negative of OSIL row~$t$.
Let $\lambda\in\R^T$ satisfy $\lambda_{T-1}=0$ and $\lambda_t<\frac14$ for $1\le t\le T-1$, and put $A_t=h(1-4\lambda_t)$ and $\bar y_t(\lambda)=(\lambda_t-\lambda_{t-1})/(2A_t)$.
Then every exactly feasible point $x$ satisfies
\begin{equation}\label{eq:dtoc5-squares}
  f(x)=q(\lambda)+h\sum_{t=0}^{T-1}\Bigl(u_t+\frac{\lambda_t}2\Bigr)^2+\sum_{t=1}^{T-1}A_t\bigl(y_t-\bar y_t(\lambda)\bigr)^2\ \ge\ q(\lambda),
\end{equation}
where $q(\lambda)=h(1-4\lambda_0)-\lambda_0-\frac h4\sum_{t=0}^{T-1}\lambda_t^2-\sum_{t=1}^{T-1}(\lambda_{t-1}-\lambda_t)^2/(4A_t)$.
\end{proposition}

\begin{proof}
At a feasible point $f(x)=f(x)+\sum_t\lambda_tc_t(x)$; collect the terms by variable.
The control $u_t$ contributes $hu_t^2+h\lambda_tu_t=h(u_t+\lambda_t/2)^2-h\lambda_t^2/4$.
The state $y_t$ with $1\le t\le T-1$ contributes $A_ty_t^2+(\lambda_{t-1}-\lambda_t)y_t=A_t(y_t-\bar y_t)^2-(\lambda_{t-1}-\lambda_t)^2/(4A_t)$, where $A_t>0$.
The fixed state $y_0=1$ contributes $h-\lambda_0-4h\lambda_0$, and $y_T$ contributes $\lambda_{T-1}y_T=0$.
\end{proof}

The identity needs no state enclosure or variable bound. It is an affine split with the dynamics rows kept in the bags and slopes $-\lambda_t$ on the separators $y_{t+1}$; $\lambda_{T-1}=0$ is necessary because the free $y_T$ would otherwise make the Lagrangian unbounded below.

\begin{theorem}[\inst{dtoc5} bracket]\label{thm:dtoc5-bracket}
Let $\hat x=(\hat u,\hat y)$ be the stored exactly feasible rational point (\cref{sec:points}), and set $\hat\lambda_t=-2\hat u_t$ for $t\le T-2$ and $\hat\lambda_{T-1}=0$.
Then $\hat\lambda_t<0$ for $t\le T-2$, and
\begin{align*}
  \vstar&\ge q(\hat\lambda)\ge5.38967211918114046742396649913627186883131268,\\
  \vstar&\le f(\hat x)\le5.38967211918114046742396649913627186883131341,
\end{align*}
with $f(\hat x)-q(\hat\lambda)\le\sci{7.21}{-43}$.
\end{theorem}

\begin{proof}
All rows hold at $\hat x$ in rational arithmetic, and $\hat u_t>0$ for $t\le T-2$, so $\hat\lambda_t<0<\frac14$ and \cref{prop:dtoc5-identity} applies; the certificate evaluates $q(\hat\lambda)$ and $f(\hat x)$.
\end{proof}

In the certificate (\cref{app:dtoc5-proof}), two separately written exact codes \arith{E} both give the displays; its evidence level is \evid{hand}+\evid{stored}.
The lower end in \cref{thm:dtoc5-bracket} is a rigorous rational evaluation below the dual function at $\hat\lambda$; we claim neither an exact dual value nor a zero duality gap.
The optimal value is attained (\cref{cor:dtoc5-attained}), and the Lagrangian dual and the order-one semidefinite relaxation are tight to within \sci{7.21}{-43} (\cref{rem:dtoc5-sdp}).

\subsection{Richer split forms}\label{sec:split-richer}

\inst{optcdeg2} controls a forced oscillator with a quadratic velocity term over $N=50{,}000$ Euler steps of length $h=1/2500$.
With $a=0.02h$ and $b=0.2h$ the model is
\[
  \min\ \frac h2\sum_{t=0}^Ny_t^2\quad\text{s.t.}\quad y_{t+1}=y_t+hv_t,\quad v_{t+1}=v_t+hu_t-ay_t-bv_t^2\quad(0\le t<N),
\]
with $|u_t|\le\frac15$, $v_t\ge-1$ for $0<t<N$, $y_0=10$, $v_0=v_N=0$ and $y_1,\dots,y_N$ free; the CUTEst problem OPTCDEG2 has $b/4$ in place of $b$ \citep{luksan2026-cutest-source-forms-for-dtoc5}.
At our exactly feasible point the control is $+\frac15$ for $3092\le t<47290$, $-\frac15$ before and after, and fractional at the two switching stages.
With the discrete costates $(p^y_t,p^v_t)$ as slopes, the stage residual of the affine split contains $-b\,p^v_{t+1}v_t^2$, which is concave in $v_t$ where $p^v>0$, that is, on both $u=-\frac15$ arcs.
On these arcs our point does not minimize the stage residual, so this affine split is not exact; a floating-point evaluation with refined costates puts its bound about 0.6258 below the optimal value.
The certificate therefore adds curvature in the velocity only: its split function on the state $(y_t,v_t)$ is $\varphi_{t-1}(y,v)=p^y_t\,y+p^v_t\,v+\frac{q_t}{2}(v-\bar v_t)^2$, where $\bar v$ is the velocity of a refined local solution.
The curvature $q_t$ makes the stage residual slightly convex in $v$ at the refined trajectory $\bar v$ on the $u=-\frac15$ arcs (\cref{fig:optcdeg2-calibration}); validity does not depend on these choices.

\begin{theorem}[\inst{optcdeg2} bound]\label{thm:optcdeg2-bound}
Let the split functions $\varphi_t$ be given by the stored binary64 split data $p^y,p^v,q,\bar v$, read as exact rationals.
Every exactly feasible point of \inst{optcdeg2} has objective at least $B$, the rigorous rational evaluation of the bound of \cref{lem:optcdeg2-krotov}, and $B\ge293.87607509587509237$.
An exactly feasible point has objective at most $293.87607509587509328$, so $\gapabs\le\sci{8.98}{-16}$ and $\gaprel\le\sci{3.06}{-18}$.
\end{theorem}

\Cref{app:optcdeg2-proof} proves it with \cref{lem:split-bound}(a) and velocity enclosures $V_t\subseteq[-1,1.2135]$ from one forward and one backward pass through the dynamics with outward rounding.
In each stage residual, $y$ enters as a quadratic with positive leading coefficient and is eliminated in closed form, and $u$ enters through one affine expression; the remaining polynomials in $v$, of degree at most four, are bounded below with exact Sturm sequences, and the stage bounds, rounded down to the grid $2^{-200}$, are summed as integers.
In the certificate (\cref{app:optcdeg2-proof}), one exact code \arith{E} gives the displayed bound, and a separately written binary64 interval code certifies the weaker bound 293.8760750958728; its evidence level is \evid{stored}.

\inst{chain}$N$ is the trapezoidal transcription, with $N$ intervals, of the hanging chain of COPS~2.0 \citep{dolan2001-benchmarking-optimization-software-with-cops}: a chain of length~4 hangs between $(0,1)$ and $(1,3)$ with minimal potential energy (\cref{app:chain-model}).
The variables are the heights $x_i$ and slopes $u_i$ ($0\le i\le N$); the dynamics and the length row are equalities, and only the end heights are fixed.

\begin{theorem}[\inst{chain} bounds]\label{thm:chain-bound}
For $N\in\{50,100,200,400\}$, the optimal value of \inst{chain}$N$ is at least a binary64 number $L_N$, displayed rounded down in \cref{tab:closures}.
Exactly feasible points with coordinates in a real quadratic field give $\gapabs\le\sci{9.58}{-15}$, \sci{1.01}{-14}, \sci{9.34}{-15} and \sci{9.78}{-15}, respectively.
\end{theorem}

\Cref{app:chain} gives the proof.
The linear rows express the variables $x_i$ and $u_i$ through the vertex heights $z_1,\dots,z_N$ of a polyline over a fixed horizontal mesh; its two end pieces depend only on the end values $z_1$ and $z_N$ (\cref{lem:chain-polyline}).
Summation by parts then writes the objective with the cumulative length $v$ of the polyline as a second state, started at a free value $V$.
For a parameter $H>0$ let $G_H(v)=\frac12\bigl(v\sqrt{H^2+v^2}+H^2\operatorname{asinh}(v/H)\bigr)$ and $c_H=2G_H(1/(2N))$.
For the split $\varphi_k(z,v)=G_H(v)-vz-k\,c_H$, a closed-form inequality makes every interior stage residual nonnegative, with equality exactly on discrete catenaries (\cref{lem:chain-calibration}).
So the objective is at least an explicit function $B_N(z_1,z_N;V,H)$ of the end values, for all $V$ and all $H>0$ (\cref{prop:chain-endvalue}).
In this split the multiplier of each height difference is the cumulative length, which varies with the configuration as the tension of a real chain does; a fixed multiplier on the length row, a reverse-convex equality, gives only about 4.77 against an optimum of about 5.07 (floating-point grid estimate).
A two-dimensional interval branch and bound over the bounded region of \cref{lem:chain-window} then proves, on each leaf box, either $B_N\ge L_N$ for some $(V,H)$ chosen for that box, or that the box contains no feasible end values (\cref{lem:split-bound}(c), with (a) in each box).
In the certificate (\cref{app:chain-bnb}), two separately written interval codes \arith{I} certify the same four values; its evidence level is \evid{rerun}, because the search trees are not stored.

\inst{catmix}$N$ transcribes the catalyst mixing problem of COPS~2.0 \citep{dolan2001-benchmarking-optimization-software-with-cops}.
Its rows are $P(u_{i+1})x_{i+1}=Q(u_i)x_i$ ($0\le i<N$), with controls $u_i\in[0,1]$, states $x_i\in\R^2$, $x_0=(1,0)$, and $2\times2$ matrices that are affine in the control and contain a coefficient $c$ (\cref{app:catmix-model}).
The objective $x_{1,N}+x_{2,N}-1$ is minus the yield, about $-0.048$.
The GAMS file has $c=9a$ with $a=1/(2N)$; the OSIL file stores $c$ as a binary64 product that exceeds $9a$ by at most $5\cdot10^{-18}$.

\begin{theorem}[\inst{catmix} bounds]\label{thm:catmix-bound}
For $N\in\{100,200,400,800\}$, the optimal value of \inst{catmix}$N$ is at least a binary64 number $L_N$, displayed rounded down in \cref{tab:closures}; this also holds for the model with $c=9a$.
Exactly feasible points give $\gapabs\le\sci{1.85}{-13}$, \sci{1.90}{-11}, \sci{6.81}{-11} and \sci{1.49}{-10}, respectively.
\end{theorem}

\Cref{app:chain} also gives this proof: for fixed controls the rows determine the states, and the stage maps $M(u)=Q(u)P(u)^{-1}$ are entrywise nonnegative, so the objective plus one is $\mathbf 1^\top P(u_N)^{-1}M(u_{N-1})\cdots M(u_1)Q(u_0)x_0$, with each control in one factor.
The value functions of this dynamic program are concave, positively homogeneous and nonnegative on $\R^2_+$, and the certificate computes chord minorants of them backward over the stages on grids of several thousand rays (\cref{lem:split-minorant}(b)).
At fixed controls the objective does not increase when $c$ increases, which transports the bounds to $c=9a$.
In the certificate (\cref{app:catmix-computation}), one code gives the displayed bounds in arithmetic \arith{F} with $+,-,\times,\div$ only, and a separately written code certifies bounds at most \sci{3.7}{-10} lower.
The bounds have evidence level \evid{rerun}, since the per-stage chord values are not stored; the points have level \evid{stored} (exact rational state evaluation).

\subsection{Affine split plus window: \inst{lukvle10}}\label{sec:split-window}

\inst{lukvle10} is problem 5.10 of Luk\v{s}an and Vl\v{c}ek with $n=1000$, in the form distributed with CUTEst \citep{gould2015-cutest-a-constrained-and-unconstrained,luksan2026-cutest-source-forms-for-dtoc5}.
Its 1000 variables are free, and with $F(a,b)=(a^2)^{b^2+1}+(b^2)^{a^2+1}$ it reads
\[
  \min\ \sum_{i=0}^{499}F(x_{2i},x_{2i+1})\quad\text{s.t.}\quad 1-x_j+3x_{j+1}-2x_{j+2}-2x_{j+1}^2=0\quad(0\le j\le997).
\]
Row $j$ defines $x_{j+2}$ from $(x_j,x_{j+1})$, so the feasible points are indexed by $(x_0,x_1)$.

\begin{theorem}[\inst{lukvle10} bracket]\label{thm:lukvle10-bracket}
The optimal value of \inst{lukvle10} satisfies
\[
  352.2380254050784\le\vstar\le352.2380254064957;
\]
hence $\gapabs\le\sci{1.42}{-9}$ and $\gaprel\le\sci{4.03}{-12}$.
\end{theorem}

\Cref{app:lukvle10-proof} gives the proof.
Dualizing rows 0 to 993 and keeping rows 994 to 997 (\cref{lem:split-bound}(d), with $Y$ the set where the kept rows hold) separates the Lagrangian into 497 two-variable pair problems and a window over the last three pairs, the row-form analogue of \cref{prop:split-window}.
The window covers the end of the chain, where the best known point leaves the saddle fixed point $-1/\sqrt2$ of the recurrence and the pair Lagrangians are nonconvex.
The dualized rows give each squared variable a coefficient $q\in[-0.8156,0]$, so each pair Lagrangian grows at least like $(1+q)t^2$ in each variable~$t$, and each infimum over $\R^2$ is a minimum over an explicit box (\cref{lem:lukvle10-coercive}).
Least-squares KKT multipliers at MINLPLib's point p5, with one common value on the middle of the chain, leave 35 distinct problems for interval branch and bound.
The summed branch-and-bound tolerances dominate the certified gap; global optimality of the exactly feasible point is not proved.
In the certificate (\cref{app:lukvle10-proof}), one code gives the displayed bound and a separately written one certifies 352.238025369202; both use interval arithmetic \arith{I} with \texttt{mpmath}'s interval $\exp$ and $\log$, and the evidence level is \evid{rerun}.

\subsection{Large stages: \inst{waterno2}}\label{sec:split-waterno2}

The \inst{waterno2} instances schedule nine pumps in four stations of a network with one reservoir and three tanks \citep{huang2019-optimal-operation-of-water-supply} over $T\in\{6,9,12,18,24\}$ hourly periods, at minimal pumping cost.
Removing the $3(T-1)$ rows that copy the tank levels between periods and one horizon row leaves $T$ periods of 166 variables (9 binary) and 203 rows each; the separators are the three tank levels.
The horizon row requires the total inflow to be at least $r_T$, the exact sum of the demands, so it is equivalent to a terminal-volume row on the last period (\cref{lem:waterno2-horizon}).

For level slopes $\lambda_t\in\R^3$ and a horizon multiplier $\mu\ge0$, the period Lagrangian gives $\vstar\ge\mu r_T+\sum_tB_t$ for lower bounds $B_t$ on the Lagrangian of each period (\cref{prop:waterno2-lagrangian}).
This is \cref{lem:split-bound}(a) for an affine split, with the horizon row dualized as well, which is valid because $\mu(r_T-\sum_th_t)\le0$ at every feasible point ($h_t$ is the inflow of period~$t$).
It certifies \inst{waterno2_09}--\inst{waterno2_24}, with multipliers from a bundle method that was stopped early.
Because pump costs depend nonconvexly on the tank levels, an affine split can leave the stage minimizers of consecutive periods at different level vectors, which weakens the bound.
For \inst{waterno2_06} we therefore divide the range of the three tank levels at each separator into 148 to 240 cells, each with its own slope vector (\cref{prop:split-cellwise}); the bound rises from 263.735099 ($\gaprel\le7.27\%$) to 278.230573 ($\gaprel\le1.68\%$).
The arc weights $b_t$ follow, through an exact linear correction (\cref{lem:waterno2-reuse}), from 49,315 rigorous bounds for single periods over pairs of level boxes, and the exact shortest-path value is $39157472136693483/2^{47}$.
Each period or pair bound comes from a spatial branch and bound with polyhedral relaxations and safe LP bounds \citep{neumaier2004-safe-bounds-in-linear-and}, which discards every box whose bound reaches a preset target value.
SCIP and HiGHS supply only multipliers, targets and dual vectors, which affect the strength of the bounds but not their validity; this matters because SCIP~10 returned wrong optimal values on some period subproblems of \inst{waterno2_06} (\cref{sec:solvers-invalid}).

\begin{theorem}[\inst{waterno2} dual bounds]\label{thm:waterno2-bounds}
For $T=6,9,12,18,24$, the optimal value of \inst{waterno2_T} is at least 278.230573, 824.834692, 2089.754565, 4790.820715 and 6576.151388, respectively, and exactly feasible points give $\gaprel\le1.68\%$, $10.82\%$, $6.90\%$, $4.87\%$ and $5.90\%$.
\end{theorem}

The bounds are 1.68 to 6.21 times the best listed dual bounds (\cref{tab:unclosed}); no instance is closed.
Regenerating the SCIP-derived targets gives slightly lower, also valid bounds for \inst{waterno2_18} and \inst{waterno2_24} (\cref{app:repro-regen}).
The point for \inst{waterno2_06} is MINLPLib's listed point p4 made exactly feasible (\cref{sec:points}); the other four are ours.
All five points use identities of the decimal data such as $0.7^3=0.343$; this identity matters for the SCIP error of \cref{sec:solvers-invalid}.
In the certificate (\cref{app:waterno2-verification}), either of two separately written branch-and-bound codes yields every value, one with node bounds in \arith{F} and one in \arith{E}; corrections, sums and the shortest path are exact \arith{E}.
Its evidence level is \evid{rerun}, because the search trees are not stored.
